\documentclass[10pt,a4paper]{amsart}
\usepackage[margin=19mm]{geometry}
\usepackage{amsmath,amssymb,mathtools}
\usepackage[scr=rsfs,cal=euler]{mathalfa}
\usepackage{xfrac}
\usepackage[unicode,hidelinks,bookmarks=false]{hyperref}
\newcommand{\ie}{i.e.\ }
\newcommand{\cf}{cf.\ }
\newcommand{\resp}{respectively\ }
\newcommand{\Subsection}{\subsection}
\providecommand{\nobreakdash}{\nobreak\hbox{-}}
\setlength{\emergencystretch}{2em}
\multlinegap0pt
\usepackage{amssymb}
\usepackage{tikz}
\newtheorem{prop}{Proposition}
\newcommand{\A}{{\mathcal A}}
\title{Free line arrangements with low maximal multiplicity}
\author{Alexandru Dimca, Lukas K\"uhne, Piotr Pokora}
\date{Source: arXiv:2505.01733v4 (CC BY 4.0)}
\begin{document}
\maketitle

\noindent\emph{Source and licence.} Free line arrangements with low maximal multiplicity. Version arXiv:2505.01733v4 (CC BY 4.0), distributed by arXiv under the Creative Commons Attribution 4.0 International licence, http://creativecommons.org/licenses/by/4.0/, as recorded in the arXiv licence metadata for this exact version and shown on the abstract page https://arxiv.org/abs/2505.01733v4. Full licence notice: \texttt{licenses/arxiv-2505.01733-CC-BY-4.0.txt}.\par\smallskip

\noindent\textit{Editorial scope.} Counted unit: the article's prop propCOMB (source0.tex lines 583--589). The article's complete original proof of it is delivered with it (source0.tex lines 590--600, one \texttt{proof} environment of 11 lines). No mathematics is written, repaired or re-derived: the statement and the proof are the article's own lines, in the article's own notation, and no retained line is substituted or re-flowed.  No further source-local context is needed: every cross-reference of every retained line resolves inside the two delivered spans.  Nothing was omitted from any retained span, so the package carries no omission record.  The retained lines cite 1 works, whose entries are transcribed verbatim from the article's own bibliography source in the e-print and delivered with short editorial labels recorded in the item's reference mapping.  No retained line contains a figure environment, an image-inclusion command or drawing code, so no image is delivered and the package declares no asset dependency.  Editorial formatting only, in addition: an AMS \texttt{amsart} preamble that restates the article's own declarations and macros used by the retained lines, namely the environments \texttt{prop} on separate counters, together with the standard packages those lines need.  The retained environments print in this document as Proposition 1 in order of appearance, whereas the article prints the same items with the numbering of its own full document.  The complete native source of the article as distributed in the arXiv e-print for this exact version is bundled unchanged as \texttt{sources/177526/source0.tex}.
\begin{prop}[cf. {\cite[Proposition 6.5]{Osaka}}]
\label{propCOMB}
Let $\A \subset \mathbb{P}^{2}$ be a free arrangement of $d$ lines. Then
$$\sum_{r\geq 2}(r-1)n_{r} \leq \bigg\lfloor \frac{(d-1)(d+3)}{4}\bigg\rfloor,$$
and this bound is sharp, for instance it is achieved by $\mathcal{A}_{13}$.

\end{prop}

\begin{proof}
Recall that the freeness of $\A$ implies that 
$$d_{1}d_{2} = d_{1}(d-d_{1}-1) = \sum_{r\geq 2}(r-1)n_{r}-d+1,$$
which gives us 
$$-d_{1}^{2}+d_{1}(d-1) - \bigg(\sum_{r\geq 2}(r-1)n_{r}-d+1\bigg)=0.$$
The above condition implies that the discriminant satisfies
$$\triangle_{d_{1}} = (d-1)^2 - 4\bigg(\sum_{r\geq 2}(r-1)n_{r}-d+1\bigg) \geq 0,$$
and this gives us
$$\sum_{r\geq 2}(r-1)n_{r} \leq \bigg\lfloor \frac{(d-1)(d+3)}{4}\bigg\rfloor.$$
Then it is easy to check that for $\A_{13}$ we do indeed get equality.
\end{proof}

\begin{thebibliography}{99}
\bibitem[Osaka]{Osaka} A. Dimca, D. Ibadula, D. A. Măcinic, Numerical invariants and moduli spaces for line arrangements. \textit{Osaka J. Math.} \textbf{57(4)}: 847 -- 870 (2020).
\end{thebibliography}
\end{document}
